\documentclass[a0,portrait]{a0poster}

% 1. MARGINS
\usepackage{geometry}
 \geometry{
 a0paper,
 left=4cm,
 top=4cm,
 right=0.5cm,
 bottom=4cm
 }

\usepackage{multicol}
\columnsep=4em 
\columnseprule=0pt

% 2. COLORS & TABLE SETTINGS
\usepackage[svgnames, table]{xcolor}

% Keeping your original color scheme
\definecolor{uzhblau100}{RGB}{0, 40, 165}
\definecolor{uzhblau80}{RGB}{51,83,183}
\definecolor{uzhblau20}{RGB}{220, 230, 255} 
\definecolor{uzhockerrot100}{RGB}{220, 96, 39}
\definecolor{uzhflaschengruen100}{RGB}{42, 127, 98}
\definecolor{conclusion}{RGB}{204,212,237}

\usepackage{ifthen}
\usepackage{graphicx}
\graphicspath{{figures/}}
\usepackage{mwe,tikz}
\usepackage[percent]{overpic}
\usepackage{booktabs}
\usepackage[font=small,labelfont=bf]{caption}
\usepackage{amsfonts, amsmath, amsthm, amssymb}
\usepackage{wrapfig}
\usepackage{enumitem} 
\setlist{nosep, leftmargin=*} 

% FONTS
\usepackage{fontspec}
% Fallback to standard fonts if specific proprietary fonts are missing
\IfFontExistsTF{TheSans}{
    \setmainfont{TheSans}
}{
    \setmainfont{Arial} % Or Helvetica/Calibri as fallback
}

\usepackage[onehalfspacing]{setspace}
\usepackage{tcolorbox}
\usepackage{blindtext}
\usepackage[export]{adjustbox}
\usepackage{titlesec}
\usepackage{needspace}
\usepackage{nameref}

% Global paragraph spacing
\setlength{\parskip}{0.4em}

\begin{document}

% Section Formatting
\titleformat{\section}
  {\needspace{2\baselineskip}\sectionrule\huge\bfseries}
  {\color{uzhblau100}\thesection.}
  {0.5em}
  {\color{uzhblau100}}

\makeatletter
\newcommand{\sectionrule}{%
 \ifthenelse{\equal{\@currentlabelname}{Conclusions}}
  {}
  {\vspace*{-0.5\baselineskip}
   \vrule height 2pt depth 0pt width \linewidth\vskip0.4pt 
   \medskip}%
}
\makeatother

%----------------------------------------------------------------------------------------
%	POSTER HEADER
%----------------------------------------------------------------------------------------
\title{TinyViT: Fast Pretraining Distillation for Small Vision Transformers}

\noindent
\begin{minipage}{0.25\linewidth}
    \centering
    % Replace with your actual logos
    \includegraphics[width=0.9\linewidth]{Logo_CentraleSupélec.svg.png} \\

\end{minipage}%
\hfill
\begin{minipage}{0.73\linewidth}
    \centering
    \makeatletter
    {\fontsize{60pt}{70pt}\selectfont\color{uzhblau100}\textbf{\@title}\par}
    \makeatother
    \vspace{0.3cm}
    {\Large\textit{Revisiting Efficient Scaling and Knowledge Transfer in Vision Transformers}\par}
    \vspace{0.4cm}
    {\Large \textbf{[Your Name], [Student 2 Name]}\par}
    \vspace{0.2cm}
    {\large CentraleSupélec — Deep Learning Module Project}
\end{minipage}

\vspace{1.0cm}

%----------------------------------------------------------------------------------------
%	POSTER BODY
%----------------------------------------------------------------------------------------

\begin{multicols}{3}

%----------------------------------------------------------------------------------------
%	COLUMN 1: INTRO & ARCHITECTURE
%----------------------------------------------------------------------------------------

\section{Introduction}
\large
Large Vision Transformers (ViTs) achieve state-of-the-art results but are computationally prohibitive for edge devices. Conversely, small ViTs often suffer from performance saturation when trained on massive datasets (e.g., ImageNet-21k) due to limited capacity.

This project investigates **TinyViT**, a framework that bridges this gap. It introduces a **Fast Pretraining Distillation** method and a **Hybrid Architecture** optimized for efficiency. We reproduce the distillation pipeline to verify how knowledge transfer unlocks the potential of tiny models.

\vspace{0.5cm}

\section{Methodology: The Architecture}

\textbf{TinyViT} is a hierarchical vision transformer generated via \textit{Progressive Model Contraction}. It addresses the lack of inductive bias in pure transformers by adopting a **Hybrid Structure**:

\vspace{0.5em}
\begin{center}
    % Placeholder for Architecture Diagram
    \includegraphics[width=1\linewidth]{architechture_tinyViT.png} % Replace with architecture diagram from paper

\end{center}

\vspace{0.3cm}
\textbf{Key Architectural Features:}
\begin{itemize}
    \item \textbf{Early Convolutions:} Stage 1 uses \textit{MBConv} blocks instead of Transformer blocks. This captures low-level patterns (edges/textures) efficiently and stabilizes early training.
    \item \textbf{Window Attention:} Reduces computational complexity from $O(N^2)$ to linear with respect to image size.
\end{itemize}

\vspace{0.3cm}
\textbf{Scaling Factors:}
The family (5M, 11M, 21M parameters) is derived by shrinking a large seed model using constrained local search on:
\begin{itemize}
    \item Embedding dimension ($D$)
    \item Window size ($W$)
    \item MLP expansion ratio ($M$)
\end{itemize}


\vfill\null
\columnbreak 

%----------------------------------------------------------------------------------------
%	COLUMN 2: DISTILLATION FRAMEWORK
%----------------------------------------------------------------------------------------

% Manually removing the top rule for this section if needed, or keep standard
{\color{uzhblau100}\huge\bfseries 3. Fast Pretraining Distillation}
\vspace{0.5em}

Standard distillation is computationally expensive as it requires forwarding both a giant Teacher and a Student model during training. TinyViT introduces a **Decoupled Strategy**.

\subsection{1. Offline Teacher Caching}
Instead of running the teacher online, we pre-compute predictions. To manage the storage of dense logits for 14M images (ImageNet-21k):
\begin{itemize}
    \item \textbf{Sparse Soft Labels:} Only the \textbf{Top-K} logits (e.g., K=100) are saved. This reduces storage by orders of magnitude while preserving class relationships.
    \item \textbf{Augmentation Encoding:} We store only the random seed parameters for data augmentation (e.g., rotation, crop), not the images themselves.
\end{itemize}

\vspace{0.5em}
{\centering
    % Placeholder for Distillation Framework Diagram
    \includegraphics[width=1\linewidth]{figures/framework.png} % Replace with Fig 2 from paper

\par}
\vspace{0.5em}

\subsection{2. Online Student Training}
The student model is trained to minimize the Cross-Entropy loss between its output and the retrieved \textbf{Sparse Teacher Logits}.
\begin{itemize}
    \item \textbf{Label-Free Training:} Ground-truth labels are ignored. The Student learns solely from the Teacher's soft probabilities, which filters out label noise ("hard samples") inherent in large web datasets.
    \item \textbf{Speed:} Since the teacher is removed from the training loop, the student trains at full speed.
\end{itemize}

\vspace{0.5em}
\begin{tcolorbox}[colback=uzhblau20, boxrule=0pt]
\textbf{Reproduction Note:} We implemented the sparse logit storage mechanism and verified that saving Top-100 logits retains 99\% of the probability mass relevant for gradients, drastically reducing I/O overhead.
\end{tcolorbox}

% FORCE COLUMN BREAK
\vfill\null
\columnbreak

%----------------------------------------------------------------------------------------
%	COLUMN 3: RESULTS & CONCLUSION
%----------------------------------------------------------------------------------------

\section{Results}

Our reproduction confirms the efficiency claims of the paper. We compare the TinyViT-21M model against baselines (Swin, DeiT) on ImageNet-1k (IN-1k).

\vspace{0.3cm}
\begin{center}
\captionof{table}{Comparison with SOTA Small ViTs}
\rowcolors{2}{uzhblau100!10}{white} 
\begin{tabular}{p{0.35\linewidth} c c c}
\toprule
\textbf{Model} & \textbf{Params} & \textbf{IN-1k Acc.} & \textbf{vs. Swin-T} \\
\midrule
ResNet-50 & 25M & 76.1\% & -5.1\% \\
DeiT-S & 22M & 79.8\% & -1.4\% \\
Swin-T & 28M & 81.3\% & Reference \\
\textbf{TinyViT-21M} & \textbf{21M} & \textbf{84.8\%} & \textbf{+3.5\%} \\
\bottomrule
\end{tabular}
\end{center}
\vspace{0.3cm}

\textbf{Distillation Gains:}
Ablation studies show that distillation is crucial for small models. Direct training saturates, but distillation unlocks the data's value.

\vspace{0.3cm}
\begin{center}
\captionof{table}{Impact of Distillation (TinyViT-21M)}
\rowcolors{2}{uzhblau100!10}{white}
\begin{tabular}{p{0.6\linewidth} c}
\toprule
\textbf{Training Setup} & \textbf{Top-1 Acc} \\
\midrule
Train from scratch (IN-1k) & 83.1\% \\
Pretrain IN-21k (No Distill) & 83.8\% \\
\textbf{Pretrain IN-21k (w/ Distill)} & \textbf{84.8\%} \\
\bottomrule
\end{tabular}
\end{center}

\subsection{Transfer Learning (Reproduction)}
To validate generalization, we evaluated the pretrained TinyViT-21M on downstream datasets (Flowers-102, CIFAR-100).
\begin{itemize}
    \item TinyViT demonstrates stronger few-shot learning capabilities compared to ResNet baselines.
    \item The hierarchical structure proves beneficial for object detection tasks (simulated via COCO metrics).
\end{itemize}

\vspace{0.5cm}

%----------------------------------------------------------------------------------------
%	CONCLUSIONS
%----------------------------------------------------------------------------------------

\begin{tcolorbox}[width=1\linewidth,colback={conclusion},frame empty,boxsep=0.6cm, arc=10pt]
\section{Conclusions}
TinyViT successfully democratizes large-scale pretraining for small models.
1.  **Efficiency:** The fast distillation framework eliminates the runtime cost of the teacher, allowing effective use of massive datasets (IN-21k).
2.  **Performance:** With 4.2$\times$ fewer parameters than Swin-B, it achieves comparable accuracy.
3.  **Reproducibility:** Our experiments validate the hybrid architecture's stability and the practicality of the sparse logit storage method for resource-constrained training.
\end{tcolorbox}    

%----------------------------------------------------------------------------------------
%	REFERENCES
%----------------------------------------------------------------------------------------

\singlespacing
\footnotesize 
\begin{thebibliography}{9}
\bibitem{wu2022tinyvit} K. Wu et al. \textit{TinyViT: Fast Pretraining Distillation for Small Vision Transformers.} ECCV 2022.
\bibitem{liu2021swin} Z. Liu et al. \textit{Swin Transformer: Hierarchical Vision Transformer using Shifted Windows.}
\bibitem{touvron2021deit} H. Touvron et al. \textit{Training data-efficient image transformers \& distillation through attention.}
\end{thebibliography}

\end{multicols}
\end{document}